\subsection{\texorpdfstring{THE MAPS \((\mathbf{G} / \mathbf{T})\times \mathbf{T}\to \mathbf{G}\) AND \((\mathbf{G} / \mathbf{T})\times \mathbf{A}\rightarrow \mathbf{G}_r\)}{THE MAPS (\textbackslash mathbf\{G\} / \textbackslash mathbf\{T\})\textbackslash times \textbackslash mathbf\{T\}\textbackslash to \textbackslash mathbf\{G\} AND (\textbackslash mathbf\{G\} / \textbackslash mathbf\{T\})\textbackslash times \textbackslash mathbf\{A\}\textbackslash rightarrow \textbackslash mathbf\{G\}\_r}}

The map \((g,t)\mapsto gtg^{-1}\) from G × T to G induces by passage to the quotient a morphism of analytic manifolds

\[
f: (\mathrm{G} / \mathrm{T}) \times \mathrm{T} \to \mathrm{G},
\]

which is surjective (§2, no. 2, Th. 2). By restriction, \(f\) induces a surjective morphism

\[
f _ {r}: (\mathrm{G} / \mathrm{T}) \times \mathrm{T} _ {r} \rightarrow \mathrm{G} _ {r}.
\]

By composition with \(Id_{G/T} \times exp_{T}\) , we obtain morphisms, also surjective,

\[
\begin{array}{c} \varphi : (\mathrm{G/T}) \times \mathfrak {t} \to \mathrm{G}, \\ \varphi_ {r}: (\mathrm{G/T}) \times \mathfrak {t} _ {r} \to \mathrm{G} _ {r}; \end{array}
\]

finally, if A is an alcove of t, \(\varphi_{r}\) induces a surjective morphism

\[
\varphi_ {\mathrm{A}}: (\mathrm{G/T}) \times \mathrm{A} \to \mathrm{G} _ {r}.
\]

We define a right operation of W on G/T as follows: let \(w \in W\) and \(u \in G/T\) ; lift w to an element n of \(\mathrm{N}_{\mathrm{G}}(\mathrm{T})\) and u to an element g of G. Then the image of gn in G/T does not depend on the choice of n and g; we denote it by u.w.

For this operation, W operates freely on G/T: indeed, with the preceding notations, assume that u.w = u; then \(gn \in gT\) , so \(n \in T\) and w = 1.

We define a right operation of W on \((\mathrm{G}/\mathrm{T}) \times \mathrm{T}\) by

\[
(u, t). w = (u. w, w ^ {- 1} (t)), \quad u \in \mathrm{G} / \mathrm{T}, t \in \mathrm{T}, w \in \mathrm{W}
\]

and a right operation of \(W_{a}^{\prime}\) on (G/T) × t by

\[
(u, x). \omega = (u. \overline {{\omega}}, \omega^ {- 1} (x)), \quad u \in \mathrm{G/T}, x \in \mathfrak {t}, \omega \in \mathrm{W} _ {a} ^ {\prime},
\]

where \(\overline{\omega}\) is the image of \(\omega\) in the quotient \(W_{a}^{\prime}/\Gamma(T)=W\) .

If A is an alcove of t, and if \(H_{A}\) is the subgroup of \(W_{a}^{\prime}\) that stabilizes A, we obtain by restriction an operation of \(H_{A}\) on \((\mathrm{G}/\mathrm{T}) \times \mathrm{A}\) .

These different operations are compatible with the morphisms \(f, \varphi\) and \(\varphi_{\mathrm{A}}\) : for \(u \in \mathrm{G} / \mathrm{T}, t \in \mathrm{T}, x \in \mathfrak{t}, y \in \mathrm{A}, w \in \mathrm{W}, \omega \in \mathrm{W}_a', h \in \mathrm{H}_{\mathrm{A}}\) , we have

\[
f ((u, t). w) = f (u, t), \varphi ((u, x). \omega) = \varphi (u, x), \varphi_ {\mathrm{A}} ((u, y). h) = \varphi_ {\mathrm{A}} (u, y).
\]

Lemma 4. Let \(g \in \mathrm{G}\) , \(t \in \mathrm{T}\) , and let \(\bar{g}\) be the image of \(g\) in \(\mathrm{G} / \mathrm{T}\) . Identify the tangent space of \(\mathrm{G} / \mathrm{T}\) (resp. T, resp. G) at \(\bar{g}\) (resp. t, resp. \(gtg^{-1}\) ) with \(\mathfrak{g} / \mathfrak{t}\) (resp. \(\mathfrak{t}\) , resp. \(\mathfrak{g}\) ) by means of the left translation \(\gamma(g)\) by \(g\) (resp. \(t\) , resp. \(gtg^{-1}\) ). The tangent linear map of \(f\) at \((\bar{g}, t)\) is then identified with the linear map \(f': (\mathfrak{g} / \mathfrak{t}) \times \mathfrak{t} \to \mathfrak{g}\) defined as follows: if \(z \in \mathfrak{g}, x \in \mathfrak{t}\) , and if \(\bar{z}\) denotes the image of \(z\) in \(\mathfrak{g} / \mathfrak{t}\) , then

\[
f ^ {\prime} (\bar {z}, x) = (\mathrm{Ad} g t ^ {- 1}) (z - (\mathrm{Ad} t) z + x).
\]

Let \(\mathrm{F}\) be the map from \(\mathrm{G} \times \mathrm{T}\) to \(\mathrm{T}\) such that \(\mathrm{F}(g,t) = gtg^{-1}\) . Since \(\mathrm{F} \circ (\gamma(g), \mathrm{Id}_{\mathrm{T}}) = \operatorname{Int} g \circ \mathrm{F}\) , we have \(\mathrm{T}_{(g,t)}(\mathrm{F})(gz, tx) = \mathrm{T}_t(\operatorname{Int} g) \circ \mathrm{T}_{(e,t)}(\mathrm{F})(z, tx)\) ; by Chap. III, §3, no. 12, Prop. 46,

\[
\mathrm{T} _ {(e, t)} (\mathrm{F}) (z, t x) = t ((\mathrm{Ad} t ^ {- 1}) z - z) + t x = t ((\mathrm{Ad} t ^ {- 1}) (z - (\mathrm{Ad} t) z + x))
\]

and consequently

\[
\mathrm{T} _ {(g, t)} (\mathrm{F}) (g z, t x) = g t g ^ {- 1} ((\mathrm{Ad}   g t ^ {- 1}) (z - (\mathrm{Ad}   t) z + x)).
\]

The lemma follows immediately from this formula by passage to the quotient.

PROPOSITION 4. a) Let \(g \in \mathbf{G}, t \in \mathrm{T}, x \in \mathfrak{t}\) , and let \(\bar{g}\) be the image of \(g\) in \(\mathrm{G} / \mathrm{T}\) . The following conditions are equivalent:

\begin{enumerate}
\def\labelenumi{(\roman{enumi})}
\tightlist
\item
  \(t \in \mathrm{T}_r\) (resp. \(x \in \mathfrak{t}_r\) ).
\end{enumerate}

(i bis) The element \(f(\bar{g}, t)\) (resp. \(\varphi(\bar{g}, x)\) ) is regular in G.

\begin{enumerate}
\def\labelenumi{(\roman{enumi})}
\setcounter{enumi}{1}
\tightlist
\item
  The map \(f\) (resp. \(\varphi\) ) is a submersion at the point \((\bar{g}, t)\) (resp. \((\bar{g}, x)\) ).
\end{enumerate}

(ii bis) The map \(f\) (resp. \(\varphi\) ) is étale at the point \((\bar{g}, t)\) (resp. \((\bar{g}, x)\) ).

\begin{enumerate}
\def\labelenumi{\alph{enumi})}
\setcounter{enumi}{1}
\item
  The map \(f_{r}\) (resp. \(\varphi_{r}\) , resp. \(\varphi_{A}\) ) makes \((\mathrm{G}/\mathrm{T}) \times \mathrm{T}_{r}\) (resp. \((\mathrm{G}/\mathrm{T}) \times \mathfrak{t}_{r}\) , resp. \((\mathrm{G}/\mathrm{T}) \times \mathrm{A}\) ) a principal covering of \(G_{r}\) with group W (resp. \(W_{a}^{\prime}\) , resp. \(H_{A}\) ).
\item
  The equivalence of (i) and (i bis) is clear; that of (ii) and (ii bis) follows from the relations \(\dim((\mathrm{G}/\mathrm{T})\times\mathrm{T})=\dim((\mathrm{G}/\mathrm{T})\times\mathfrak{t})=\dim(\mathrm{G})\) . By Lemma 4, f is a submersion at the point \((\bar{g},t)\) if and only if \(g=t+Im(Ad t-Id)\) , which means that t is regular. Finally, since \(\varphi=f\circ(Id_{G/T}\times\exp_{T})\) , \(\varphi\) is étale at the point \((\bar{g},x)\) if and only if f is étale at the point \((\bar{g},\exp x)\) , which by the preceding means that x belongs to \(t_{r}\) .
\item
  The morphisms \(f_{r}, \varphi_{r}, \varphi_{A}\) are thus étale. On the other hand, W operates freely on G/T, and a fortiori on \((\mathrm{G}/\mathrm{T}) \times \mathrm{T}\) . Let \(g, g'\) in G and \(t, t'\) in \(T_{r}\) be such that \(f(\bar{g}, t) = f(\bar{g}', t')\) ; then \(\operatorname{Int} g^{-1} g'\) maps \(t'\) to t, and hence normalizes T, since \(\mathrm{T} = \mathrm{Z}(t)_{0} = \mathrm{Z}(t')_{0}\) , and the class w of \(g^{-1} g'\) in W maps \((\bar{g}, t)\) to \((\bar{g}', t')\) . It follows that \(f_{r}\) is a principal covering with group W; this immediately implies that \(\varphi_{r}\) is a principal covering with group \(W_{a}'\) , and hence by restriction to the connected component \((\mathrm{G}/\mathrm{T}) \times \mathrm{A}\) of \((\mathrm{G}/\mathrm{T}) \times \mathbf{t}_{r}\) , that \(\varphi_{A}\) is a principal covering with group \(H_{A}\) .
\end{enumerate}

Remarks. 1) By Prop. 3 of §2, no. 4, the manifold \((\mathrm{G}/\mathrm{T}) \times \mathrm{A}\) is simply-connected. It follows that \(\varphi_{A}\) is a universal covering of \(G_{r}\) ; since \(\pi_{1}(G_{r})\) is canonically isomorphic to \(\pi_{1}(G)\) (no. 1, Prop. 1 and General Topology,

Chap. XI, in preparation), we recover the fact that \(\pi_{1}(\mathrm{G})\) can be identified with \(H_{A}\) (that is with \(\Gamma(\mathrm{T})/\mathrm{N}(\mathrm{G},\mathrm{T})\) ).

\begin{enumerate}
\def\labelenumi{\arabic{enumi})}
\setcounter{enumi}{1}
\item
  The restriction of \(\varphi_{A}\) to \(W \times A \subset (G/T) \times A\) makes \(W \times A\) a principal covering of \(T_{r}\) with group \(H_{A}\) . We thus recover Cor. 1 of Prop. 2 of no. 2.
\item
  Denote by \(\mathfrak{g}_r\) the inverse image of \(\mathrm{G}_r\) under the exponential map and by \(\varepsilon : \mathfrak{g}_r \to \mathrm{G}_r\) the map induced by \(\exp_{\mathrm{G}}\) . The map \((g, x) \mapsto (\operatorname{Ad} g)(x)\) from \(\mathrm{G} \times \mathfrak{t}_r\) to \(\mathfrak{g}_r\) defines by passage to the quotient a map \(\psi_r : (\mathrm{G} / \mathrm{T}) \times \mathfrak{t}_r \to \mathfrak{g}_r\) . We have \(\varepsilon \circ \psi_r = \varphi_r\) . Let \(w \in \mathrm{W}, \gamma \in \Gamma(\mathrm{T})\) and \(\omega \in \mathrm{W}_a'\) be such that \(\omega(z) = w(z) + \gamma\) for all \(z \in \mathfrak{t}\) ; then \(\psi_r((\bar{g}, x)\omega) = \psi_r(\bar{g}, x) - (\operatorname{Ad} g)(\gamma)\) for \(g \in \mathrm{G}, x \in \mathfrak{t}_r\) , so \(\psi_r((\bar{g}, x)\omega) = \psi_r(\bar{g}, x)\) if and only if \(\gamma = 0\) . It follows (cf.~General Topology, Chap. XI, in preparation) that \(\psi_r\) is a principal covering of \(\mathfrak{g}_r\) with group W, and that \(\varepsilon : \mathfrak{g}_r \to \mathrm{G}_r\) is a covering associated to the principal covering \(\varphi_r\) , with fibre isomorphic to the \(\mathrm{W}_a'\) -set \(\mathrm{W}_a'/\mathrm{W}\) .
\end{enumerate}

